\section{Two further quantization axes: moment and depth}
\label{sec:quant-axes}

The previous sections vary the \emph{calibration} of the kernel at a fixed format. Two further degrees of
freedom drive the outcome: \emph{when} one quantizes, and \emph{how many bits} one keeps.

\subsection{Moment: before, after, or both}

Three placements are possible for the same EWC head. \emph{Before} training (QAT, fake-quant in the loop);
\emph{after} (PTQ on an already trained FP32 model); and \emph{both}, that is a QAT run whose weights are
then exported to the calibrated integer kernel --- only this last path is faithful to the actual deployment.

Under emulation, all three moments preserve AUROC on Monitoring (\num{0.974477} before, \num{0.974649}
after, \num{0.974663} both, against \num{0.97464} in FP32): at this model scale no moment is disqualifying.
The question is settled \emph{on-board}. The "both" path was flashed and streamed: it yields \num{0.9213}
F1 on Monitoring and \num{0.9072} on Pronostia, against respectively \num{0.9173} and \num{0.8995} for the
calibrated PTQ alone (Section~\ref{sec:recovery-board}).

The gap, \num{0.004} and \num{0.0077}, is real but small. The honest reading is therefore that
\textbf{kernel calibration recovers what matters, and that QAT adds no decisive gain} on this head: when
the training pipeline cannot be modified, a properly calibrated PTQ suffices. This contradicts the common
reading whereby a PTQ collapse would necessarily call for retraining.

\begin{figure}[t]
\centering
\includegraphics[width=0.9\linewidth]{fig6_quant_moment.png}
\caption{Quantization moment. Left, the three moments \emph{PC-emulated} (hatched); right, the
\emph{on-board measured} A/B between calibrated PTQ alone and the QAT$\to$PTQ path.}
\label{fig:moment}
\end{figure}

\subsection{Depth: how far below 8 bits}

Below 8 bits the question becomes: at which depth does the metric break down, and does the granularity of
the scale factor move that limit? The emulated sweep answers both.

The limit depends on the dataset. On Monitoring, degradation stays contained down to binary
($\Delta$AUROC $=-\,\num{0.0117}$). On Pronostia, binary breaks down ($-\,\num{0.0275}$) but ternary holds
($-\,\num{0.0153}$): the usable frontier is therefore \emph{ternary}, binary being defensible only on the
easier dataset.

Granularity, in turn, pushes the breakdown back. At 2 bits on Pronostia, moving from a per-tensor to a
per-channel scale takes $\Delta$AUROC from $-\,\num{0.0464}$ to $-\,\num{0.0086}$. A negative result
deserves reporting too: the affine \emph{zero-point} buys nothing --- at 2 bits on Monitoring it even
degrades markedly ($-\,\num{0.0619}$ against $-\,\num{0.0069}$ symmetric).

\paragraph{The RAM gain only exists packed.} The crucial point is that the RAM advertised by these schemes
is \emph{theoretical} as long as weights stay in an \texttt{int8} container. Measured on-board, the
\texttt{.bss} of a 4-, 2- or 1-bit \emph{unpacked} model is strictly identical to the INT8 one: the gain
only appears with genuine bit-packing, where it grows as bits drop (\num{336}, \num{504} then
\num{572}~bytes saved on Monitoring). Unpacking is paid in latency, about \SI{55}{\micro\second} (from
\SI{67}{\micro\second} unpacked to \SI{123}{\micro\second} packed on Monitoring) --- without threatening the
Gap~2 budget. Parity against the emulator stays \num{1.000} across the twelve measured cells: the scheme is
ported losslessly.

\begin{figure}[t]
\centering
\includegraphics[width=0.9\linewidth]{fig7_quant_depth_packing.png}
\caption{Sub-INT8 depth. Left, $\Delta$AUROC by depth and granularity (\emph{PC-emulated}); right, the
\texttt{.bss} bytes actually saved by bit-packing (\emph{measured on-board}).}
\label{fig:depth}
\end{figure}
